% Cote 119: offprints, reports, courses, CV and letters (typed and printed part): archive cote 119, site page 210 (PDF page 211).
% typescript. Draft compared: none.
% AI-transcribed by Claude Opus 5.5 (claude-opus-5-5, adaptive thinking, effort high), Message Batches API, batch r119_full (msgbatch_01DwZNm5Y7YsKBWyAN4sFs8w), 2026-11-13; not yet checked by eye.
\scanpage{119/210}
\mspage{- 15 -}

LISTE DES TRAVAUX CITES,SUSCITES OU INFLUENCES PAR LES TRAVAUX DE A. GROTHENDIECK

M. ARTIN, Algebraization of Formal Moduli, I (in Global Analysis, pp. 21-71,
University of Tokyo Press, Princeton University Press, 1968), II Exis-%
tence of modifications,Annals of Mathematics, Vol. 91, pp. 88-135
(1970).

M.F. ATIYAH et F. HIRZEBRUCH, Riemann-Roch theorems for differentiable manifolds,
Bull.Amer. Math. Soc., vol. 65, pp. 276-281 (1959).

P. BERTHELOT, Cohomologie cristalline des schémas propres et lisses de caracté-%
ristique $p > 0$ , Thèse, Université Paris VII, 1971 (paraîtra dans
Lecture Notes of Math. chez Springer).

P. BERTHELOT et L. ILLUSIE, Classes de Chern en cohomologie cristalline, C.R. Acad.
Sci.Paris t. 270, pp. 1695-1697 (22 juin 1970) et p. 1750-1752
(29 juin 1970).

A. BOREL et J.P. SERRE, Le théorème de Riemann-Roch (d'après des résultats \unsure{iné|ine}-%
dits de A. Grothendieck), Bull. Soc. Math. France, t. 86,
pp. 97-136 (1958).

P. DELIGNE, Théorie de Hodge I (Actes du Congrès International des mathématiciens,
Nice 1970) et II, Publications Math. n° 40, pp. 5-57 (1971).

P. DELIGNE et D. MUMFORD, The irreducibility of the space of curves of a given
genus, Pub. Math. n° 36, pp. 75-110 (1969).

M. DEMAZURE, Motifs des variétés algébriques, Sém. Bourbaki n° 365, 1969/70.

A. DOUADY, Le Problème des modules pour les sous-espaces analytiques compacts...,
Ann. Inst. Fourier, vol. 16, pp. 1-98 (1966).

O. FORSTER et K. KNORR, Relativ-analytische \unsure{Räume|Raume} und die \unsure{Kohärenz|Koharenz} von Bild-%
garben,Inventiones Math. Vol. 16, pp. 113-160 (1972).

I.M. GELFAND et N. Ja. VILENKIN, Les distributions, tome 4, Applications de
l'analyse harmonique, Dunod, Paris, 1968 (traduction).

J. GIRAUD, Cohomologie non abélienne, Grundlehren des Math. Wiss. Bd. 179, 1971,
Springer.

\margin{210}
